\begin{proposition}[obstruction to a specified repair]\hypertarget{tag-0307}{}\label{push:src:B.3.3}\label{push:B33}
\StatementRecord{push:src:B.3.3}{proved}{0307}{\SeriesRef{push:src:B.3.1}, \SeriesRef{push:src:A.1.2}}

A cyclic cusp-root construction using a
fractional translation by the two cusp
periods cannot repair \ref{push:src:B.3.1} while retaining
the reduced Mumford fibre for $r\ne0$.
More precisely, let $q=s^d$, $u=r^n$,
and suppose a cyclic deck generator has
translation $d^{-1}(aG_1+bG_2)$ in the
additive notation for the two periods.
Assume its $r$-component shift is integral.
Then its $r$-component shift is
$nb/d\bmod n$, whereas the multiplicity
of every generic quotient $q$-component
is a positive multiple of
\[
 \frac{d}{\gcd(d,a,b)}.
\]
A nonzero $r$-component shift therefore
forces a nonreduced $q$-fibre. In
particular a shift $-1\bmod12$ in this
class of constructions forces multiplicity
at least twelve.
\end{proposition}


\begin{proof}
\begin{enumerate}[label=\arabic*.,leftmargin=*]
\item In the coordinates of \ref{push:src:A.1.2} the two
periods have $q$-valuation vectors
$(1,0)$ and $(0,1)$; their $r$-valuation
vectors are $(0,0)$ and $(0,n)$.
After $q=s^d$, the period lattice on
the $q$-central component indices is
$d\mathbb Z^2$. A regular toroidal
resolution has indices in
$(\mathbb Z/d)^2$. The proposed deck
generator translates them by $(a,b)$.
The order of that translation is
$d/\gcd(d,a,b)$, so the subgroup
preserving a component has order
$h=\gcd(d,a,b)$. Its subgroup fixing
a generic point has some order
$h_0$ dividing $h$; unit factors may
make $h_0$ smaller. That isotropy
acts on the normal coordinate $s$
by a primitive $h_0$th root of
unity. The quotient normal
coordinate is $s^{h_0}$, and
$q=(s^{h_0})^{d/h_0}$ has
multiplicity $d/h_0$, a positive
multiple of $d/h$. Unit factors
cannot decrease this lower bound.

\item Reducedness would require $h_0=d$,
and therefore $h=d$,
so both $a,b$ are divisible by $d$.
Then $nb/d$ is zero modulo $n$.
If $nb/d$ is relatively prime to
$n$, its order in $\mathbb Z/n$
is $n$, which divides the order
$d/h$ of the full component
translation; this gives the final
lower bound. The residue calculations
and examples for the requested
exponents are independently recorded
in \texttt{ab\_push\_cusp\_root\_repair.py}.
This rules out the specified repair,
not a cusp model with a different
geometry or a different comparison
to the punctured family.
\end{enumerate}
\end{proof}
